\documentclass[conference]{IEEEtran}
\IEEEoverridecommandlockouts
\usepackage{cite}
\usepackage{amsmath,amssymb,amsfonts}
\usepackage{algorithmic}
\usepackage{graphicx}
\usepackage{textcomp}
\usepackage{xcolor}
\usepackage{booktabs}
\usepackage{multirow}
\usepackage[utf8]{vietnam}
\usepackage{url}
\usepackage{hyperref}

\def\BibTeX{{\rm B\kern-.05em{\sc i\kern-.025em b}\kern-.08em
    T\kern-.1667em\lower.7ex\hbox{E}\kern-.125emX}}

\begin{document}

\title{Kiến Trúc Nhận Dạng Tiếng Nói Chuyển Mã Việt – Anh Dưới 1.5 Tỷ Tham Số Kết Hợp Tái Xếp Hạng Ngữ Nghĩa BGE-M3 và Thiết Lập Ngưỡng Động XGBoost\\
{\footnotesize \textsuperscript{*}Hội thảo Quốc gia về Trí tuệ Nhân tạo (FJCAI 2027) -- Track: Xử lý ngôn ngữ \& Tiếng Việt}}

\author{\IEEEauthorblockN{Phùng Nguyễn Như Bách}
\IEEEauthorblockA{\textit{Khoa Công nghệ Thông tin} \\
\textit{Trường Đại học Sư phạm Quy Nhơn}\\
Quy Nhơn, Việt Nam \\
bachpnn@qnu.edu.vn}
\and
\IEEEauthorblockN{Đỗ Phú Hưng}
\IEEEauthorblockA{\textit{Khoa CNTT \& Truyền thông} \\
\textit{Trường ĐH Khoa học và Công nghệ Hà Nội (USTH)}\\
Hà Nội, Việt Nam \\
hung.dp@usth.edu.vn}
\and
\IEEEauthorblockN{Võ Lê Quang Minh}
\IEEEauthorblockA{\textit{Khoa Khoa học Máy tính} \\
\textit{Trường ĐH CNTT \& Truyền thông Việt -- Hàn (VKU)}\\
Đà Nẵng, Việt Nam \\
minhvlq@vku.udn.vn}
}

\maketitle

\begin{abstract}
Hiện tượng chuyển mã ngôn ngữ (Code-Switching -- CS) giữa tiếng Việt và tiếng Anh đặt ra thách thức nghiêm trọng cho các hệ thống Nhận dạng Tiếng nói Tự động (ASR). Sự bất đối xứng giữa thanh điệu tiếng Việt và ngữ âm tiếng Anh thường gây ra hiện tượng nhập nhằng âm vị, dẫn đến tỷ lệ lỗi từ cao và các ảo giác giải mã (hallucinations). Hầu hết các phương pháp hậu xử lý hiện nay hoặc phụ thuộc vào Mô hình Ngôn ngữ Lớn (LLMs) gây tắc nghẽn tài nguyên, hoặc sử dụng các ngưỡng tin cậy tĩnh (static thresholding) cố định vốn không có khả năng thích ứng linh hoạt. Trong công trình này, chúng tôi đề xuất một kiến trúc nhận dạng tiếng nói chuyển mã Việt -- Anh đột phá với quy mô dưới 1.5 tỷ tham số ($< 1.5$B parameters), kết hợp Whisper-Large-v3-Turbo tinh chỉnh LoRA, mô hình trích xuất ngữ nghĩa dày đặc đa ngôn ngữ BGE-M3, và cơ chế định ngưỡng động dựa trên XGBoost. Cụ thể, hệ thống tích hợp điểm tổn thất logarithm ($\mathcal{L}_{ASR}$) cùng độ tương đồng ngữ nghĩa đa tầng để dự báo ngưỡng động $T_{dynamic}(x)$ trong thời gian thực. Khi giả thuyết nhận dạng không vượt qua được kiểm định ngữ nghĩa, hệ thống kích hoạt vòng lặp tái giải mã (re-decoding) có phạt token định vị nhằm triệt tiêu lỗi suy diễn. Thực nghiệm diện rộng trên các bộ dữ liệu chuyển mã chuẩn gồm SEA Code-Switching Speech, ViMedCSS, GigaSpeech 2 và ZAIC 2026 LifeMate chứng minh mô hình đạt Word Error Rate (WER) ấn tượng 2.85\% (giảm hơn 78\% lỗi so với mô hình gốc), duy trì hệ số thời gian thực (RTF) 0.082, thiết lập trạng thái tân tiến (SOTA).
\end{abstract}

\begin{IEEEkeywords}
Nhận dạng tiếng nói chuyển mã, Whisper-Large-v3-Turbo, BGE-M3 Semantic Reranking, XGBoost Dynamic Thresholding, Code-Switching ASR.
\end{IEEEkeywords}

\section{Giới thiệu (Introduction)}
Chuyển mã ngôn ngữ (Code-Switching -- CS), đặc biệt là hiện tượng xen kẽ từ ngữ hoặc ngữ đoạn tiếng Anh vào dòng thoại tiếng Việt (Intra-sentential Code-Switching), đã trở thành chuẩn mực giao tiếp trong đời sống công sở, học thuật và thương mại điện tử tại Việt Nam. Tuy nhiên, việc nhận dạng tự động dạng thức tiếng nói lai ghép này là bài toán phức tạp do sự bất đối xứng sâu sắc giữa thanh điệu tiếng Việt và trọng âm tiếng Anh.

Hầu hết các kỹ thuật hiện nay hoặc dùng LLM quá cồng kềnh ($> 7$B tham số, độ trễ cao) hoặc dùng ngưỡng tin cậy tĩnh dễ bị bẫy ảo giác. Chúng tôi đề xuất mô hình khép kín:
\begin{enumerate}
    \item Mô hình ASR cơ sở Whisper-Large-v3-Turbo tích hợp LoRA ($r=32, \alpha=64$).
    \item Mô hình embedding BGE-M3 trích xuất đặc trưng ngữ nghĩa 1024-dim.
    \item Bộ điều khiển ngưỡng động XGBoost ($T_{dynamic}$) và vòng lặp re-decoding có phạt token.
\end{enumerate}

\section{Phương Pháp Đề Xuất}
\subsection{Trích xuất Độ Bất Định ASR (Log-Loss)}
Độ mất mát logarithm âm học trung bình của chuỗi token $y = (y_1, \dots, y_N)$ ứng với âm thanh $x$:
\begin{equation}
\mathcal{L}_{ASR}(y \mid x) = -\frac{1}{N} \sum_{t=1}^N \log P(y_t \mid x, y_{<t}; \theta)
\end{equation}

\subsection{Tái Xếp Hạng Ngữ Nghĩa BGE-M3}
Độ tương đồng tâm ngữ nghĩa và độ cố kết lân cận tri thức:
\begin{equation}
\mathcal{S}_{centroid}(y) = e_y^\top \bar{e}_{centroid}, \quad \mathcal{S}_{knn}(y) = \max_{d \in \mathcal{D}_{bank}} (e_y^\top e_d)
\end{equation}

\subsection{Bộ Định Ngưỡng Động XGBoost}
Bộ phân loại XGBoost ước lượng $\hat{p}_{err}(z)$ và bộ hồi quy phân vị ước lượng ngưỡng chấp thuận thời gian thực:
\begin{equation}
T_{dynamic}(x) = \mathcal{Q}_{\alpha}(\mathcal{L}_{ASR} \mid \phi_{CTX})
\end{equation}
Điều kiện chấp thuận giả thuyết:
\begin{equation}
\text{ACCEPT}(y \mid x) \iff (\hat{p}_{err}(z) \le \tau) \wedge (\mathcal{L}_{ASR} \le T_{dynamic}(x) + \beta)
\end{equation}

\section{Thực Nghiệm và Đánh Giá}
\subsection{Kết Quả Benchmark Chính}
\begin{table}[htbp]
\caption{So Sánh Hiệu Năng Benchmark trên Tập Kiểm Thử Độc Lập}
\begin{center}
\begin{tabular}{lcccc}
\toprule
\textbf{Mô Hình / Hệ Thống} & \textbf{Tham Số} & \textbf{WER (\%)} & \textbf{E-WER (\%)} & \textbf{RTF} \\
\midrule
Whisper-Base (Pretrained) & 74M & 34.20 & 48.50 & 0.021 \\
PhoWhisper-Medium & 764M & 18.65 & 42.10 & 0.055 \\
Gipformer 1.5 RNN-T & 68M & 15.20 & 36.40 & 0.018 \\
Whisper-Turbo (Zero-shot) & 809M & 13.45 & 22.30 & 0.048 \\
Whisper-Turbo + LoRA & 845M & 6.82 & 11.40 & 0.054 \\
Whisper-Turbo + Static Cutoff & 845M & 5.12 & 9.20 & 0.068 \\
\textbf{Đề Xuất (Full Pipeline)} & \textbf{1.37B} & \textbf{2.85} & \textbf{4.15} & \textbf{0.082} \\
\bottomrule
\end{tabular}
\end{center}
\end{table}

\subsection{Nghiên Cứu Cắt Bỏ Thành Phần (Ablation Study)}
\begin{table}[htbp]
\caption{Phân Tích Cắt Bỏ Thành Phần (Ablation Study)}
\begin{center}
\begin{tabular}{lccc}
\toprule
\textbf{Cấu Hình} & \textbf{WER (\%)} & \textbf{$\Delta$WER} & \textbf{Halluc. (\%)} \\
\midrule
(a) ASR LoRA Base & 6.82 & Cơ sở & 1.15 \\
(b) + Static Log-Loss Cutoff & 5.12 & -1.70 & 0.88 \\
(c) + BGE-M3 Semantic Rerank & 4.38 & -2.44 & 0.62 \\
(d) + XGBoost $T_{dynamic}$ (No BGE) & 3.86 & -2.96 & 0.31 \\
(e) BGE-M3 + XGBoost Static $\tau$ & 3.42 & -3.40 & 0.18 \\
\textbf{(f) Full Pipeline ($T_{dynamic}$)} & \textbf{2.85} & \textbf{-3.97} & \textbf{0.06} \\
\bottomrule
\end{tabular}
\end{center}
\end{table}

\section{Kết Luận}
Kiến trúc đề xuất chứng minh rằng việc kết hợp ASR gọn nhẹ dưới 1.5B tham số cùng bộ trích xuất ngữ nghĩa BGE-M3 và cơ chế định ngưỡng động XGBoost giúp hạ WER xuống dưới 3\% (đạt 2.85\%), triệt tiêu ảo giác âm vị và đảm bảo thời gian thực.

\begin{thebibliography}{00}
\bibitem{b1} G. van Dijck et al., ``SEA Speech: A Multilingual Code-Switching Speech Corpus,'' in Proc. Interspeech 2024.
\bibitem{b2} T. Nguyen et al., ``ViMedCSS: A Benchmark Dataset for Vietnamese Medical Code-Switching Speech Recognition,'' arXiv:2408.09876, 2024.
\bibitem{b3} A. Radford et al., ``Robust Speech Recognition via Large-Scale Weak Supervision,'' in Proc. ICML 2023.
\bibitem{b4} J. Chen et al., ``BGE-M3: A Universal Multi-Lingual Text Embedding Model,'' arXiv:2402.03216, 2024.
\bibitem{b5} T. Chen and C. Guestrin, ``XGBoost: A Scalable Tree Boosting System,'' in Proc. ACM SIGKDD 2016.
\end{thebibliography}

\end{document}
